\chapter{Preliminarii}
\label{ch:preliminarii}

\section{Proiectarea și Abordarea Cercetării}
\label{sec:two-axes}

Pentru a răspunde la întrebarea formulată mai sus, în loc să optimizez o singură arhitectură pentru un singur set de date, am realizat o comparație de-a lungul a două axe simultan.

Prima axă este \textbf{paradigma}: cele șapte modele evaluate (\S\ref{sec:paradigms}) acoperă patru moduri structural distincte de a reprezenta cum arată \enquote{normalul} -- o \textbf{graniță} trasată într-un spațiu indus de un kernel (One-Class SVM), o \textbf{partiționare} a spațiului de caracteristici prin diviziuni aleatoare (Isolation Forest), un \textbf{bottleneck de compresie} învățat, care trebuie să renunțe la orice, cu excepția structurii dominante a datelor normale (cele trei autoencodere bazate pe reconstrucție), și un model \textbf{predictiv} al evoluției pas cu pas a datelor normale (LSTM predictiv, Transformer cauzal). Fiecare paradigmă codifică o presupunere implicită diferită despre ce anume face un punct anormal, ceea ce ridică întrebarea centrală a acestei lucrări: care presupunere se potrivește cărui tip de anomalie.

A doua axă este \textbf{setul de date}: Credit Card Fraud, Yahoo Webscope S5 și CIC-IDS2017 (\S\ref{sec:datasets-overview}) acoperă un spectru de structură temporală -- de la un set pur tabelar, fără ordonare semnificativă între rânduri, trecând prin serii de timp cu adevărat secvențiale, până la fluxuri de rețea care admit atât o perspectivă plată, cât și una temporală asupra acelorași evenimente. Rularea fiecărei paradigme pe fiecare punct al acestui spectru face posibilă atribuirea unui rezultat dat \emph{paradigmei}.

\subsection{Ipoteze și Comportament Așteptat}
\label{sec:expected-behaviour}

Pe baza acestor două axe, anticipez două tipare în rezultate.

În primul rând, mă aștept ca paradigmele care oferă modelului o structură temporală inerentă a datelor -- prin construcție, nu doar prin faptul că le sunt oferite ferestre de intrare -- să performeze în general mai bine pe seturile de date cu adevărat secvențiale (Yahoo S5 și, într-o măsură mai mică, varianta temporală a CIC-IDS2017), comparativ cu paradigmele care tratează fiecare eșantion ca independent.

În al doilea rând, mă aștept ca diferite paradigme să fie mai bine adaptate la a detecta diferite \emph{tipuri} de anomalii, chiar în cadrul aceluiași set de date. Dacă această presupunere este corectă, atunci compunerea mai multor paradigme complementare într-un ansamblu ar trebui să ducă la o îmbunătățire a performanței față de oricare model individual, atâta timp cât modelele componente exploatează relații diferite din date, și nu doar variații ale aceleiași relații. Această a doua ipoteză este testată direct atât prin analiza pe tip de atac din \S\ref{ch:results} (în special pentru CIC-IDS2017, unde este disponibilă eticheta tipului de atac), cât și prin studiul de caz al modelului ansamblu din \S\ref{ch:explainability}.

În următoarele secțiuni vor fi prezentate detaliile algoritmilor și seturilor de date alese.

% \section{Current State of Anomaly Detection Research}
% \label{sec:sota}

% % TODO: back this paragraph with concrete citations (comparative surveys,
% % benchmark papers) once the bibliography is filled in.
% Much of the published work on anomaly detection either proposes a new architecture and benchmarks it against a handful of established baselines within a single data modality, or surveys the field at a conceptual level without actually running the same models, under the same protocol, across substantially different kinds of data. Comparative studies tend to stay within one modality at a time -- time-series anomaly detection surveys on one side, tabular fraud-detection benchmarks on the other -- so it remains comparatively rare to see how the same paradigm-level assumption (a bottleneck, a boundary, a predictive objective) fares when the underlying data changes shape entirely, from independent tabular rows to genuinely sequential signals to network flow records. This is the specific gap the present thesis addresses: rather than proposing a new model, the two sections below introduce the four established paradigms and the three datasets that, taken together, make it possible to ask that paradigm-versus-modality question directly, under a single, shared experimental protocol (\S\ref{sec:eval-protocol}).

\section{Algoritmi de Detecție a Anomaliilor}
\label{sec:paradigms}

\subsection{One-Class SVM}

One-Class SVM (formularea $\nu$-SVM) învață o graniță în jurul datelor de antrenare, într-un spațiu de caracteristici indus de un kernel. Trei hiperparametri, fixați identic pe toate cele trei seturi de date (\texttt{kernel='rbf'}, \texttt{nu=0.01}, \texttt{gamma='scale'} -- vezi \S\ref{sec:architecture}), determină împreună forma acestei granițe (detalii despre rolul fiecăruia: Anexa~\ref{sec:model-mechanisms}). La testare, distanța semnată față de graniță (pozitivă în interior, negativă în exterior) este negată astfel încât, în concordanță cu toate celelalte modele din acest proiect, \emph{scor mai mare = mai anormal}.

\subsection{Isolation Forest}

Isolation Forest construiește un ansamblu de arbori care împart fiecare recursiv spațiul de caracteristici pe baza unor caracteristici și puncte de divizare alese aleator, folosind lungimea medie a căii necesare pentru a izola un punct drept semnal de anomalie: punctele din regiuni rare de densitate scăzută (anomaliile) se izolează în mai puține diviziuni, deci ating o lungime medie a căii mai mică decât punctele benigne tipice. Scorul returnat este această lungime medie a căii, negată, pentru convenția consistentă \emph{scor mai mare = mai anormal}. Parametrul \texttt{contamination} este lăsat la valoarea implicită \texttt{'auto'} și nu are niciun efect asupra rezultatelor, întrucât datele de antrenare sunt exclusiv benigne, iar pragul de operare este ales extern prin baleierea curbei PR (\S\ref{sec:eval-protocol}) -- detalii: Anexa~\ref{sec:model-mechanisms}.

\subsection{Autoencoder Feedforward (MLP)}

Autoencoder Feedforward este bazat pe reconstrucție: o pereche simetrică encoder/decoder de straturi \texttt{Linear} + \texttt{ReLU} comprimă vectorul de intrare printr-un bottleneck de dimensionalitate redusă, apoi îl reconstruiește, fiind antrenat să minimizeze eroarea medie pătratică (MSE) de reconstrucție. Bottleneck-ul mai îngust decât intrarea forțează o reprezentare comprimată a structurii dominante \emph{benigne}, în loc de funcția identitate. La testare, scorul de anomalie este MSE per eșantion: eșantioanele benigne se reconstruiesc cu eroare mică, cele anormale -- pentru care bottleneck-ul nu a fost antrenat -- se reconstruiesc slab.

\subsection{Autoencoder LSTM}

Autoencoder LSTM aplică același principiu de reconstrucție, dar pentru secvențe: un strat LSTM encoder consumă întreaga secvență și păstrează doar starea ascunsă finală (și starea de celulă) ca vector sumar de dimensiune fixă, care este apoi repetată la fiecare pas ca intrare pentru un LSTM decoder, ale cărui ieșiri sunt proiectate înapoi la dimensionalitatea originală. Scorul de anomalie este, din nou, MSE per eșantion, calculat pe toți pașii și toate caracteristicile simultan.

\subsection{LSTM Predictiv}

LSTM Predictiv folosește același bloc recurent, dar cu un obiectiv diferit: este antrenat să prezică pasul $t+1$ din pașii $0..t$, la fiecare poziție simultan -- fără separare encoder/decoder și fără bottleneck, starea ascunsă la fiecare pas fiind proiectată direct într-o predicție prin un strat liniar. Scorul de anomalie este MSE dintre valorile prezise și cele reale ale pasului următor.

\subsection{Autoencodere Bazate pe Transformer}

Sunt folosite două variante, ambele înlocuind coloana vertebrală recurentă cu un \texttt{Trans\-form\-er\-En\-cod\-er} (blocuri de auto-atenție + feedforward, codificare pozițională sinusoidală): \textbf{Bottleneck Transformer AE}, bazat pe reconstrucție -- atenție non-cauzală urmată de mediere globală (\emph{global-average-pooling}) într-un vector comprimat printr-un bottleneck liniar (analog stării ascunse din LSTM AE), apoi extins înapoi; scor: MSE de reconstrucție -- și \textbf{Causal Transformer}, predictiv, echivalentul bazat pe atenție al LSTM Predictiv (mască de atenție cauzală, fără bottleneck; scor: MSE de predicție). Detalii mecanism: Anexa~\ref{sec:model-mechanisms}.

Ambele variante folosesc deliberat un singur strat \texttt{TransformerEncoder}, păstrând costul de antrenare/inferență și numărul de parametri la aproximativ jumătate din ce ar necesita un design encoder-decoder echivalent, fără a sacrifica abilitatea de a face fie comprimare-și-reconstrucție, fie predicție cauzală. Nu există o implementare Causal Transformer pentru setul de date credit card (\S\ref{ch:results}).

\section{Seturile de Date}
\label{sec:datasets-overview}

\subsection{Credit Card Fraud Detection (Setul de Date A)}

Setul de date Kaggle Credit Card Fraud Detection conține 284{,}807 tranzacții bancare, dintre care 492 (0.172\%) sunt frauduloase.
Cele 30 de caracteristici prezente sunt: \texttt{Time} (secunde scurse de la prima tranzacție din set), 28 de caracteristici transformate prin PCA \texttt{V1}--\texttt{V28} (anonimizate din motive de confidențialitate, deja decorelate și aproximativ zero-medie prin construcție), și \texttt{Amount} (valoarea brută a tranzacției, puternic asimetrică spre dreapta).
\texttt{Time} reprezintă secunde scurse, nu o axă temporală structurată -- nu există periodicitate, tendință sau noțiune de sesiune de exploatat. Acest lucru face din Setul de Date A cazul de control deliberat, \enquote{fără ordine temporală reală}.

Setul de date este disponibil pe Kaggle~\cite{kaggle_creditcard}, alături de un notebook exploratoriu dedicat gestionării datelor dezechilibrate~\cite{kaggle_janiobachmann}, care a servit drept inspirație pentru o parte din scripturile acestui proiect.

\subsection{Yahoo Webscope S5 (Setul de Date B)}

Yahoo Webscope S5 constă din patru familii de benchmark (A1--A4), fiecare o colecție de mai multe serii de timp univariate independente, cu anomalii etichetate, dar fiecare construită pentru a izola un tip diferit de anomalie:

\begin{center}
\resizebox{0.88\textwidth}{!}{%
\begin{tabular}{@{}lllll@{}}
\toprule
Benchmark & Serii & Origine & Tip de anomalie & Rata de anomalii \\
\midrule
A1 & 67  & Metrici reale de server & Valori punctuale aberante & 1.76\% \\
A2 & 100 & Sintetic, o singură componentă periodică + tendință & Valori punctuale aberante & 0.33\% \\
A3 & 100 & Sintetic, 3 componente de frecvență în interacțiune & Contextuală (deviație de fază) & 0.56\% \\
A4 & 100 & Sintetic, aceeași structură, fără sezonalitate & Schimbări de nivel & 0.50\% \\
\bottomrule
\end{tabular}%
}
\end{center}

A1 este real și zgomotos (cicluri zilnice/săptămânale dezordonate); A2--A4 sunt sintetice și construite special astfel încât tipul de anomalie să fie controlat independent de forma seriei de bază. Acest lucru este ceea ce permite capitolului de Rezultate să atribuie diferențele de performanță \emph{tipului de anomalie} (punctuală vs.\ contextuală vs.\ schimbare de nivel), și nu unor diferențe confundante între seturile de date.

\subsection{CIC-IDS2017 (Setul de Date C)}

CIC-IDS2017 (Canadian Institute for Cybersecurity) este un benchmark de detecție a intruziunilor în rețea, captat pe parcursul a cinci zile: $\sim$2.8 milioane de fluxuri de rețea, fiecare descris de $\sim$80 de caracteristici statistice extrase prin CICFlowMeter (durată, numărul și rata pachetelor/octeților, statistici ale timpului dintre sosiri, numărul de flag-uri TCP, dimensiuni de fereastră etc.). Luni este trafic 100\% benign; 14 tipuri distincte de atacuri (variante DoS/DDoS, PortScan, forță brută, atacuri web, Bot, Infiltration, Heartbleed) sunt injectate de marți până vineri.

Se știe că distribuția brută/Kaggle a acestui set de date este substanțial coruptă \parencite{engelen2021troubleshooting}: aproximativ 26\% din fluxuri sunt artefacte de tip \enquote{appendix} TCP, care moștenesc eticheta de atac a fluxului părinte fără a conține trafic de atac real; o mare parte din fluxurile etichetate nominal ca atac nu au încărcătură utilă reală (etichetate doar formal, nu pe bază de comportament); iar clasa \texttt{DoS Hulk} corespunde unei simulări de atac defectă, care nu se execută niciodată efectiv împotriva serverului țintă. Coloanele de identificare prezente (\texttt{Flow ID}, IP sursă/destinație, \texttt{Timestamp}) ar permite unui model să \enquote{învețe pe scurtătură} după adresa IP sau fereastra de captură. Acest proiect folosește versiunea curățată și reextrasă de la \textcite{engelen2021troubleshooting}, în locul CSV-urilor originale Kaggle/UNB, și elimină explicit coloanele de identificare (\S\ref{sec:preprocessing}).

Fluxurile CIC-IDS2017 sunt agregate statistice la nivel de sesiune. Acest lucru face din acest set de date un caz de tranziție perfect: admite o perspectivă validă de tip plat/tabelar (un flux = un eșantion, analog setului de Date A), o perspectivă temporală validă (fluxuri grupate și ordonate per gazdă, analog setului de Date B).